\documentclass[11pt]{article}

\usepackage[margin=1in]{geometry}
\usepackage[T1]{fontenc}
\usepackage[utf8]{inputenc}
\usepackage{amsmath,amssymb,mathtools}
\usepackage{array,booktabs,longtable,tabularx}
\usepackage{enumitem}
\usepackage{xcolor}
\usepackage{listings}
\usepackage[hidelinks]{hyperref}
\usepackage{microtype}

\newcommand{\ALCIRCC}{\mathcal{ALCI}_{\mathrm{RCC5}}}
\newcommand{\RCC}{\mathrm{RCC5}}
\newcommand{\EQ}{\mathsf{EQ}}
\newcommand{\PP}{\mathsf{PP}}
\newcommand{\PPI}{\mathsf{PPI}}
\newcommand{\PO}{\mathsf{PO}}
\newcommand{\DR}{\mathsf{DR}}
\newcommand{\Rel}{\mathsf{Rel}}
\newcommand{\Comp}{\operatorname{Comp}}
\newcommand{\St}{\operatorname{St}}
\newcommand{\Safe}{\operatorname{Safe}}
\newcommand{\core}{\operatorname{core}}
\newcommand{\tp}{\operatorname{tp}}
\newcommand{\sigmaG}{\sigma_g}
\newcommand{\code}[1]{\texttt{#1}}
\newcommand{\status}[1]{\textsc{#1}}
\newcommand{\sev}[1]{\textsc{#1}}

\setlist[itemize]{leftmargin=2em}
\setlist[enumerate]{leftmargin=2em}
\renewcommand{\arraystretch}{1.18}
\setlength{\LTpre}{0.5em}
\setlength{\LTpost}{0.5em}
\emergencystretch=2em

\lstdefinestyle{reviewpython}{
  basicstyle=\ttfamily\scriptsize,
  columns=fullflexible,
  keepspaces=true,
  breaklines=true,
  frame=single,
  showstringspaces=false,
  tabsize=2,
  language=Python
}
\lstdefinestyle{reviewtext}{
  basicstyle=\ttfamily\scriptsize,
  columns=fullflexible,
  keepspaces=true,
  breaklines=true,
  frame=single,
  showstringspaces=false,
  tabsize=2
}

\title{Formal Response on the Round-17 Transition-System Delta for\
\texorpdfstring{$\ALCIRCC$}{ALCI_RCC5}}
\author{GPT-5.5 Pro}
\date{13 July 2026}

\begin{document}
\maketitle

\begin{abstract}
This response tests whether the round-17 delta discharges the eleventh review's six required repairs R1--R6 for the soundness direction of the cluster-quasimodel decidability argument for $\ALCIRCC$.  The answer is negative.  Round 17 introduces a genuine finite transition-system vocabulary, but its ordered-pair transport rule is still too narrow: when a child gluing is crossed, one old endpoint may be updated by the ambient transition $T$, while the other old endpoint is a member of the parent copy and receives a birth-adjacent catalogue state rather than a $T$-image.  Definition 3.3 transports only pairs for which both endpoint transitions are $T$-defined.  This leaves a current old--old pair out of the S4 check and revives the same composition-closure failure that the work order was meant to eliminate.  A targeted finite check, WP31, witnesses the omission.  The appropriate status is therefore \emph{strongly supported, not certified}; round 17 is not acceptable as a full discharge of R1--R6 without the mixed-transport repair described below.
\end{abstract}

\tableofcontents

\section{Scope and notation}

The assessment uses the attached sources:
\begin{itemize}
  \item \code{workorder/formal\_referee\_report\_alci\_rcc5.tex}, especially the required repairs R1--R6 at lines 306--326;
  \item \code{manuscripts/cluster\_quasimodels\_round17\_transition\_system.tex}, especially Definitions 3.1--3.3 and Theorem 3.1;
  \item \code{verification/wp30\_round17\_transition\_system.py} and its attached output;
  \item the new targeted check \code{verification/wp31\_round17\_mixed\_transport\_attack.py}, included in Appendix~\ref{app:wp31}.
\end{itemize}
Line numbers below refer to \code{nl -ba} renderings of the corresponding \TeX{} source files.  The RCC5 atoms are written $\EQ,\PP,\PPI,\PO,\DR$; $\Comp$ denotes the standard RCC5 composition table; $\Rel$ denotes the full atom set.

\section{Executive verdict}

I do \emph{not} accept round 17 as discharging the eleventh review's work order.  The central defect is local and finite.  Round 17 defines:
\begin{itemize}
  \item an ambient singleton transition $T_{g\to g'}$ only for occurrences $y\notin V_k$ at a child gluing $g'$ of the copy $k$ attached at $g$ (round-17 lines 145--162);
  \item a separate clause saying that for $y\in V_k$ the state at $g'$ is read off the pattern network as a birth-adjacent catalogue state (lines 162--166);
  \item an ordered-pair transport rule (P3) that applies only when both singleton transitions are defined and then transports the pair as $(T\varsigma,T\varsigma',W)$ (lines 180--185).
\end{itemize}
Thus the written rule omits a mixed old--old pair in which one endpoint is ambient and the other endpoint is a parent-copy member.  That pair is neither created afresh nor co-patterned at the child gluing.  It must be transported, because S4 must inspect its current endpoint states and realized mutual value when assigning rows to new fresh members.  Theorem 3.1 then uses the false assertion that all pairs among pre-existing occurrences have both endpoint states updated by $T$ (lines 251--254).

\section{Part I: acceptance table for R1--R6}

\begin{longtable}{>{\raggedright\arraybackslash}p{0.08\linewidth} >{\raggedright\arraybackslash}p{0.20\linewidth} >{\raggedright\arraybackslash}p{0.64\linewidth}}
\toprule
\textbf{Item} & \textbf{Result} & \textbf{Assessment} \\
\midrule
\endfirsthead
\toprule
\textbf{Item} & \textbf{Result} & \textbf{Assessment} \\
\midrule
\endhead
\bottomrule
\endfoot
R1 & \status{Nominally discharged, but not fully} & Round 17 does introduce finite singleton states, ordered pair states, seeds, creation, transport, and a least fixed point; see round-17 lines 130--190.  Birth-adjacent states for inherited separator members are mostly computable from the catalogue because (N1) forces every pre-existing member of the new copy into the declared separator, while gluing coherence forces pattern values to agree with pre-existing values.  However, the singleton update vocabulary is split: $T_{g\to g'}$ is explicitly defined only for ambient endpoints, while parent-copy endpoints receive catalogue/birth-adjacent states by a separate read-off clause.  Since (P3) is stated only in terms of $T$, the transition system is not explicit for all actual endpoint-state updates needed by pair transport.  Also, the continued/core cases overlap on core members because (N3) makes $\core\subseteq\sigma_g$; harmless only with an explicit consistency invariant. \\
\addlinespace
R2 & \status{Not discharged} & The work order's Appendix-B intent was lockstep ordered-pair transport through each gluing.  Round 17's (P3), lines 180--185, transports $(\varsigma,\varsigma',W)$ only when both singleton transitions are $T$-defined.  It covers old--old pairs whose endpoints are both ambient.  It omits an old--old pair where one endpoint is ambient and the other is a member of the parent copy, hence has a birth-adjacent child-gluing state rather than a $T$-image.  It also does not cleanly classify the case where one endpoint is separator/co-patterned and the other is ambient.  These pairs are not re-seeded and are not transported by the stated rule. \\
\addlinespace
R3 & \status{Not discharged} & Theorem 3.1 is the right inclusion theorem, but its proof fails in case (c).  Lines 251--254 state that pairs among pre-existing occurrences have unchanged mutual value and both endpoint states update by $T$.  That is false for old parent-copy endpoints, whose next states are birth-adjacent catalogue states under Definition 3.2, lines 162--166.  The compressed sentence at lines 231--234, ``or a later gluing elsewhere in the tree order; the argument is the same with $T$ composed along the corresponding path,'' is also asserted rather than defined.  For pairs born at branch meetings or transported across sibling/descendant moves, the corresponding path includes endpoint-role changes not covered by $T$-composition.  Case (a)'s use of (N1) is otherwise sound for genuinely ambient endpoints: if $y\notin V_k$ and $s\in B_g$, the $y,s$ value is a steering value, not a pattern value. \\
\addlinespace
R4 & \status{Nominally discharged, but needs syntactic tightening} & (N1)--(N3) are stated at round-17 lines 92--107 and are used for the F2/core exclusion, the three-case row decomposition, and no-overwrite.  If enforced, they exclude the old defect where an old core occurrence is co-patterned with fresh members outside the separator.  But (N1) is written as $V_k\cap U=\sigma_g$, where $U$ is the set of occurrences existing before the step; that is phrased as an unfolding/run property.  The finite certificate should instead contain a checkable port partition and injection discipline: inherited slots plus core are exactly the separator; every other non-core port is fresh; non-core occurrence identity is reused only along its declared connected trace. \\
\addlinespace
R5 & \status{Discharged in the manuscript text} & Definition 3.3 makes pair states ordered triples $(\varsigma,\varsigma',W)$ and says the converse triple is added simultaneously; see lines 169--172.  S4 is then inherited as an oriented condition.  I do not see an unordered-pair reading left in the mathematical text, although WP30 is not itself a complete test of this point. \\
\addlinespace
R6 & \status{Nominally discharged, but not literally complete} & Q2$'$ stores a total witness map $w(k,a,\exists R.D)$; see lines 280--286, and lines 288--292 state that attached copies are those designated by $w$.  This fixes the previous hidden witness-value choice.  It determines the unfolding tree up to a canonical convention that each pending demand spawns its designated child.  It still does not literally specify sibling ordering, child indexing, or whether two demands designating the same gluing edge are shared or duplicated.  That is not the present soundness hole if a fixed convention is added, but R6 asked that the tree structure itself be a function of the certificate. \\
\end{longtable}

\section{Part II: new attack-surface findings}

\subsection{Finding 17.1: missing mixed transport revives the S4 omission}

\paragraph{Severity.} \sev{High}.  This is a local soundness-proof defect in the claimed explicit transition system.

\paragraph{Location.} Round-17 Definition 3.2, lines 145--167; Definition 3.3, lines 169--190; Theorem 3.1 proof, lines 231--264.

\paragraph{Issue.}  Consider a child gluing $h$ of the copy $k$ attached at $g$.  A pre-existing endpoint at $h$ may be:
\begin{enumerate}[label=(\roman*)]
  \item outside $V_k$, hence updated by the ambient transition $T_{g\to h}$;
  \item in $V_k$ and in the child separator, hence co-patterned at $h$;
  \item in $V_k$ but not in the child separator, hence assigned a birth-adjacent child-gluing state read from the catalogue.
\end{enumerate}
Round 17's (P3) transports old pair states only for the $(i,i)$ case.  It omits mixed cases such as $(i,iii)$.  Such a pair is old--old; it is not created by (P2), and it need not be co-patterned or seeded by (P1) at $h$.  But S4 must see precisely this current pair state when selecting values from old ambient elements to fresh elements.

\paragraph{Finite witness.}  Let $h$ introduce a fresh member $b$ with separator member $s$.  Let $z\in V_k$ be a parent-copy member that is not in $\sigma_h$, and let $y\notin V_k$ be ambient.  Assume the old relations
\[
  y\;\PP\;z\;\PP\;s,
  \qquad s\;\PPI\;b.
\]
Choose steering values
\[
  \rho(y,b)=\PO,
  \qquad
  \rho(z,b)=\DR.
\]
The one-step checks that do not use the omitted current pair can pass:
\[
  \PO\in\Comp(\PP,\PPI)=\Rel,
  \qquad
  \DR\in\Comp(\PP,\PPI)=\Rel.
\]
S3 is vacuous if $b$ is the only fresh member, and S2 can be made true by choosing types with no restrictive universals or with safe domains admitting the selected values.  But S4, if applied to the current old pair $(y,z)$, rejects the choice because
\[
  \Comp(\PP,\DR)=\{\DR\},
  \qquad
  \PO\notin\{\DR\}.
\]
Equivalently, the triangle $(y,z,b)$ is not composition-closed:
\[
  \rho(y,b)\notin\Comp\bigl(\rho(y,z),\rho(z,b)\bigr).
\]
The only reason this row can escape S4 is that (P3) did not transport the pair
\[
  (\St_g(y),\St_g(z),\PP)
\]
to the current pair
\[
  (\St_h(y),\St_h(z),\PP),
\]
where $\St_h(y)=T_{g\to h}\St_g(y)$ and $\St_h(z)$ is the birth-adjacent catalogue state of the parent-copy member $z$ at $h$.

\paragraph{Machine check.}  WP31, Appendix~\ref{app:wp31}, verifies the witness against the standard finite RCC5 composition table and then models the P-rules at state-symbol level.  Its output is:
\begin{lstlisting}[style=reviewtext]
PASS: S1/S2/S3 can accept the bad y-z-b triangle when the mixed pair is omitted.
PASS: The manuscript P3 text omits (state_y_at_child_h, birth_adjacent_state_z_at_h, PP).
PASS: Adding a mixed ambient/parent transport rule restores the pair, so S4 rejects the bad row.
\end{lstlisting}
This is not a full concept-level countermodel.  It is a finite transition-system witness that the written fixed point does not include a pair state that Theorem 3.1 needs.

\paragraph{Repair.}  Replace $T$ in (P3) by a total endpoint-update operation $U_{g\to h}$ defined on every pre-existing endpoint state relevant at the next gluing:
\[
  U_{g\to h}(x)=
  \begin{cases}
    T_{g\to h}(x), & x\notin V_k,\\
    \text{the catalogue/birth-adjacent state of }x\text{ at }h, & x\in V_k.
  \end{cases}
\]
Then close the pair set under
\[
  (\varsigma,\varsigma',W)
  \longmapsto
  \bigl(U_{g\to h}(\varsigma),U_{g\to h}(\varsigma'),W\bigr),
\]
including mixed $T$/birth-adjacent cases and simultaneous converses.  The proof of Theorem 3.1 must use this $U$-transport in case (c), not the assertion that both endpoints update by $T$.

\subsection{Finding 17.2: the ``later gluing elsewhere'' path is undefined}

\paragraph{Severity.} \sev{High/medium}.  This is the global form of Finding 17.1.

\paragraph{Location.} Round-17 Theorem 3.1 proof, lines 231--234.

The proof says that a child gluing or a later gluing elsewhere in the tree order is handled by the same argument, with $T$ composed along the corresponding path.  The corresponding path is not formally defined.  For a pair whose endpoints were born in different branches, the next relevant gluing may sit below a least common ancestor, in a sibling branch, or under a descendant where one endpoint is a member of the parent copy and the other remains ambient.  Along such a path, endpoint roles change among ambient, parent-copy, separator, core, and already-expanded side-branch occurrence.  Round 17 defines only the ambient $T$ and separately says parent-copy states are read from the catalogue.  Therefore the composed-$T$ claim is not proved and, in mixed cases, is false as stated.

\subsection{Finding 17.3: WP30 is weaker than the manuscript claims}

\paragraph{Severity.} \sev{Medium}.  WP30 remains useful, but it does not test the failing case.

The attached WP30 output says:
\begin{lstlisting}[style=reviewtext]
Part A (R3 inclusion): 120 random catalogues x 10 simulations each; containment failures: 0 (PASS)
Part B (R2 load-bearing): 69/69 catalogues with transported pairs NOT covered by a birth-only fixed point (PASS -- transport is necessary, confirming F1)
Part C (S4 => closure): 100 catalogues with S4-valid canonical steering; closure failures in simulation: 0 (PASS)
        control (S4-violating steering): 0/0 show closure failures (n/a -- none sampled)
\end{lstlisting}
The round-17 manuscript states 83/83 in lines 69--72 and 317--320, while the attached output reports 69/69.  More importantly, inspection of the harness shows that it exercises chain-like unfoldings and records ordered pairs primarily among its simulated ambient occurrences.  It does not exhaustively cover branching or the mixed ambient/parent-copy state update isolated by WP31.  Part C also has no negative-control S4-violating samples in the attached run.  Thus WP30 is supportive smoke testing, not evidence that (P3) covers the Appendix-B transport obligation.

\subsection{Finding 17.4: (N1) is phrased as a run property rather than certificate syntax}

\paragraph{Severity.} \sev{Medium/low}.  This is probably repairable, but R4 asked for checkability.

Round-17 Definition 3.1 states $V_k\cap U=\sigma_g$, where $U$ is the set of occurrences existing before the step.  That statement is natural semantically, but the finite certificate should state the corresponding syntactic discipline directly.  A checkable formulation should partition every template into core ports, inherited separator slots, and fresh ports; require each gluing to map inherited slots injectively into parent-copy or core ports; allocate fresh occurrence identities for every non-inherited non-core port at each copy; and forbid non-core identity reuse outside the connected trace generated by these maps.

\subsection{Finding 17.5: core/separator overlap needs an explicit row-consistency invariant}

\paragraph{Severity.} \sev{Low}.  Because (N3) says $\core\subseteq\sigma_g$ for every gluing, a core member is both a persistent core-row member and a current separator member.  The three-case decomposition lists continued and core cases separately, but they overlap on core elements.  The values should agree by construction and by no-overwrite; the manuscript should state this as an invariant, especially if the transition system is later mechanized.

\section{Checks that did not produce a fatal finding}

\subsection{Y1: parent-plus-core normal form}

I did not find a satisfiable concept forcing a separator member from a non-parent, non-core copy.  The restriction $\sigma_g\subseteq V_{k'}\cup\core$ is consistent with the running-intersection property expected from the base split-forest/tree-decomposition presentation: if an occurrence appears in a new bag and in an older non-parent bag, it must also appear along the path, hence in the parent once the tree is rooted.  This remains an extraction/completeness obligation to cite carefully, but I do not have a counterexample to it.

\subsection{Y2: no-overwrite lemma}

Under strict (N1)--(N3), I did not find an overwrite witness.  A previously $f$-generated pair cannot later become a pattern pair unless both old endpoints are reintroduced into a new copy; (N1)/(N2), plus exact occurrence identity, are intended to forbid that outside declared connected traces.  The lemma's conclusion looks correct once the finite occurrence-identity syntax from Finding 17.4 is made explicit.

\subsection{Y3: safety signature}

The safety signature appears certificate-determined: it is a function of the endpoint type, the current gluing's fresh-member types, and the finite catalogue.  It need not be an independent state component, provided states are indexed by the relevant gluing alphabet and the update rule specifies which gluing's fresh-member list is used.

\subsection{Y4: strong equality, folds, and width accounting}

I found no new strong-equality failure.  Core rows are persistent and pattern networks use $\EQ$ only on exact occurrence identity.  Fold-lattice bounds are no longer load-bearing for round-17 soundness.  Storing the witness map $w$ adds finite certificate data and branching instructions, but it does not by itself enlarge local pattern width; it affects scheduling conventions rather than bag size.

\section{Part III: optional round-18 material}

Part III is not reached as a positive development stage, because Parts I--II find a high-severity transition-system defect.  The next round should first repair (P3), restate Theorem 3.1 using the total endpoint-update operation $U$, and rerun a harness that includes mixed ambient/parent-copy transport and genuinely branching unfoldings.  Only after that repair would it be useful to spend proof budget on F6 width-budget derivations or W2$'$ uniformization.

\section{Required round-18 repair package}

A minimal acceptable repair package would include the following.
\begin{enumerate}[label=\textbf{M\arabic*.}]
  \item Define a total endpoint-update function or relation $U_{g\to h}$ for every pre-existing endpoint crossing from a gluing $g$ to a child gluing $h$, including ambient, parent-copy, separator, and core endpoints.
  \item Replace (P3) by ordered-pair transport under $U$ while preserving $W$, with simultaneous converse closure.
  \item Prove that $U$ is finite and effectively computable from the catalogue and gluing declaration, including inherited separator members whose rows mix pattern values and pre-existing values.
  \item Replace Theorem 3.1(c)'s ``both endpoint states update by $T$'' sentence by a case analysis over $U$.
  \item Define the path semantics for ``later gluing elsewhere'' via least common ancestors or an equivalent canonical traversal, and prove pair-state inclusion along composed $U$-updates.
  \item Strengthen WP30 or add WP32 so that it samples or exhaustively enumerates mixed ambient/parent-copy transport, separator/ambient transport, and branching-tree unfoldings.  The harness should include a negative S4 control.
\end{enumerate}

\section{Conclusion}

Round 17 is a substantial improvement over round 16: it names the right objects, fixes ordered-pair orientation in the text, stores witness choices, and isolates the open completeness obligations.  But the work order's core demand was not merely to name a finite fixed point; it was to guarantee that every actual current old--old pair state used by S4 is in that fixed point.  The stated (P3) does not do that.  The soundness argument remains \emph{strongly supported, not certified}, and round 17 should be rejected until the mixed endpoint-update transport is made explicit and proved.

\appendix

\section{WP31 targeted finite check}\label{app:wp31}

\subsection{Output}
\begin{lstlisting}[style=reviewtext]
PASS: S1/S2/S3 can accept the bad y-z-b triangle when the mixed pair is omitted.
PASS: The manuscript P3 text omits (state_y_at_child_h, birth_adjacent_state_z_at_h, PP).
PASS: Adding a mixed ambient/parent transport rule restores the pair, so S4 rejects the bad row.
\end{lstlisting}

\subsection{Source}
\begin{lstlisting}[style=reviewpython]
#!/usr/bin/env python3
"""
WP31 -- targeted finite attack on round-17 Definition 3.3 / Theorem 3.1.

It isolates the edge case not exercised by WP30: crossing from a copy k
into a child gluing h where one old endpoint is outside k (updated by T)
and the other old endpoint is a member of k but not of h's separator
(a birth-adjacent/catalogue state, not in the domain of T as written).

Under the manuscript's P3 text, the pair is not transported. If the
missing current pair is omitted from S4, a standard RCC5 triangle accepted
by S1/S2/S3 becomes composition-inconsistent. Adding the evident mixed
transport rule makes S4 reject it.
"""

REL = ("EQ", "PP", "PPI", "PO", "DR")
CONV = {"EQ": "EQ", "PP": "PPI", "PPI": "PP", "PO": "PO", "DR": "DR"}
COMP = {
    "EQ":  {"EQ": {"EQ"}, "PP": {"PP"}, "PPI": {"PPI"},
             "PO": {"PO"}, "DR": {"DR"}},
    "PP":  {"EQ": {"PP"}, "PP": {"PP"}, "PPI": set(REL),
             "PO": {"PP", "PO", "DR"}, "DR": {"DR"}},
    "PPI": {"EQ": {"PPI"}, "PP": {"EQ", "PP", "PPI", "PO"},
             "PPI": {"PPI"}, "PO": {"PPI", "PO"},
             "DR": {"PPI", "PO", "DR"}},
    "PO":  {"EQ": {"PO"}, "PP": {"PP", "PO"},
             "PPI": {"PPI", "PO", "DR"}, "PO": set(REL),
             "DR": {"PPI", "PO", "DR"}},
    "DR":  {"EQ": {"DR"}, "PP": {"PP", "PO", "DR"},
             "PPI": {"DR"}, "PO": {"PP", "PO", "DR"},
             "DR": set(REL)},
}
COMP = {a: {b: frozenset(v) for b, v in row.items()} for a, row in COMP.items()}


def check(condition, msg):
    if not condition:
        raise AssertionError(msg)


# Intended geometry at child gluing h:
# y PP z PP s, separator s PPI b, but steering chooses y PO b and z DR b.
yz = "PP"
zs = "PP"
ys = "PP"      # forced by y PP z PP s, and used as y's row to separator s
sb = "PPI"     # pattern edge s -> b
zb = "DR"      # steering value for parent member z toward fresh b

yb_bad = "PO"  # steering value for ambient y toward fresh b

# S1 checks through separator s. Both bad choices pass S1 because Comp(PP,PPI)=Rel.
check(yb_bad in COMP[ys][sb], "S1 should allow y PO b through s")
check(zb in COMP[zs][sb], "S1 should allow z DR b through s")

# S3 is vacuous for one fresh member; S2 can be made true by setting safe domains accordingly.
# S4, if the current pair (St_h(y), St_h(z), PP) is present, rejects y PO b.
required_by_s4 = COMP[yz][zb]
check(required_by_s4 == frozenset({"DR"}), f"Expected Comp(PP,DR)={{DR}}, got {required_by_s4}")
check(yb_bad not in required_by_s4, "S4 should reject the bad y-b value")

# Direct composition-closure failure of triangle (y,z,b): rho(y,b) must be in Comp(rho(y,z),rho(z,b)).
closure_ok = yb_bad in COMP[yz][zb]
check(not closure_ok, "The bad triangle should fail RCC5 composition closure")

# Now model just the P-rules at the state-symbol level.
# At parent gluing g, z is fresh in copy k and y is ambient, so P2 creates (Yg,Zg,PP).
Yg = "state_y_at_g"          # ambient state before attaching k
Zg = "birth_state_z_in_k"    # z's birth state at g
Yh = "state_y_at_child_h"    # T(Yg)
Zh = "birth_adjacent_state_z_at_h"  # read from catalogue; not T(Zg) in Def. 3.3

reachable_pairs = {(Yg, Zg, yz), (Zg, Yg, CONV[yz])}
T = {Yg: Yh}  # T is defined only for occurrences ambient at g; z in V_k has no T-state clause.

# Manuscript P3: transport only if both singleton transitions are defined.
for a, b, w in list(reachable_pairs):
    if a in T and b in T:
        reachable_pairs.add((T[a], T[b], w))
        reachable_pairs.add((T[b], T[a], CONV[w]))

missing_current_pair = (Yh, Zh, yz) not in reachable_pairs
check(missing_current_pair, "As written P3 unexpectedly transported the mixed ambient/parent pair")

# Evident repair: mixed transport maps the parent-copy endpoint to its birth-adjacent child-gluing state.
reachable_repaired = set(reachable_pairs)
parent_adjacent = {Zg: Zh}
for a, b, w in [(Yg, Zg, yz), (Zg, Yg, CONV[yz])]:
    if a in T and b in parent_adjacent:
        reachable_repaired.add((T[a], parent_adjacent[b], w))
    if a in parent_adjacent and b in T:
        reachable_repaired.add((parent_adjacent[a], T[b], w))

check((Yh, Zh, yz) in reachable_repaired, "Mixed transport repair should add the current pair")

print("PASS: S1/S2/S3 can accept the bad y-z-b triangle when the mixed pair is omitted.")
print("PASS: The manuscript P3 text omits (state_y_at_child_h, birth_adjacent_state_z_at_h, PP).")
print("PASS: Adding a mixed ambient/parent transport rule restores the pair, so S4 rejects the bad row.")
\end{lstlisting}

\end{document}
